% !TEX program = xelatex
% !TEX root = course_report.tex
% ============================================================
% SmartSurvey：面向空间异常检测的智能文献综述与知识图谱系统
% 项目实现报告（《计算机图形学前沿》期末项目）
% 中文 XeLaTeX 排版
% ============================================================
\documentclass[a4paper,11pt]{article}

% ---------- 字体与中文 ----------
\usepackage[UTF8]{ctex}
\usepackage[paper=a4paper, margin=1.6cm]{geometry}

% ---------- 表格 ----------
\usepackage{booktabs}
\usepackage{tabularx}
\usepackage{array}

% ---------- 插图 ----------
\usepackage{graphicx}
\graphicspath{{./figures/}}

% ---------- 数学与符号 ----------
\usepackage{amsmath}
\usepackage{amssymb}

% ---------- 超链接 ----------
\usepackage[hidelinks]{hyperref}
\urlstyle{same}

% ---------- 段落微调 ----------
\setlength{\parskip}{0.05em plus 0.02em minus 0.02em}
\usepackage{enumitem}
\setlist{nosep,leftmargin=*,topsep=0pt,partopsep=0pt}
\raggedbottom

% ============================================================
\begin{document}

% ============================================================
% 标题与作者
% ============================================================
\begin{center}
  {\LARGE\bfseries SmartSurvey：面向空间异常检测的智能文献综述与知识图谱系统}\\[0.3em]
  {\large —— 项目实现报告（《计算机图形学前沿》期末项目）}\\[0.5em]
  {\large 钟启帆 \quad PB24010467 \quad 数学科学学院，中国科学技术大学}\\[0.5em]
  {\large 开源仓库：\url{https://github.com/C6-H14/SmartSurvey} \quad 2026 年 8 月}
\end{center}

\vspace{0.3em}
\noindent\rule{\textwidth}{0.6pt}

% ============================================================
% 摘要与关键词
% ============================================================
\begin{abstract}
\noindent
面向工业与机器人场景的空间异常检测研究近年来呈爆发式增长，研究者人工梳理与对比成本急剧上升。
本文实现了一个面向空间异常检测与图形可视化的智能文献综述系统 SmartSurvey，以一组异构 3D/RGB-D
空间异常检测论文的 PDF 为输入，自动完成文献解析、空间视觉矩阵提取、证据约束审查与
中文 XeLaTeX 学术综述全文的动态排版生成。系统在架构上采用分层模块设计：以``核心章节识别 + 兜底页码切片''
混合解析保证鲁棒性，以``通用字段 + 领域字段''双层矩阵 Schema 组织对比指标，
以基于 \texttt{evidence\_quote in page\_text} 的 Inline Scope Binding 零丢纸事实校验拦截幻觉，
并以双层模板 XeLaTeX 物理排版引擎（含轻量语法栈扫描器与自愈重试）直接产出可贴入 Overleaf 的手稿。
在图形可视化层面，系统基于 PyVis 构建交互图形拓扑网络，将文献间主题关联直观呈现。
系统同时实现 OS Keyring → .env → 会话内存三级凭据降级与网关指数退避重试，保障公网与无头容器场景安全运行。
经 126 项自动化单元测试全量验证（100\% PASS）并开源交付。

\noindent\textbf{关键词：} 空间异常检测；3D/RGB-D 视觉；智能文献综述；PyVis 图形拓扑网络；XeLaTeX 动态排版；零丢纸事实校验
\end{abstract}

% ============================================================
% 1. 引言与课题背景
% ============================================================
\section{引言与课题背景}

\subsection{研究背景：空间异常检测的兴起}
空间异常检测旨在识别工业现场、自动化实验室与机器人作业环境中偏离正常模型的异常对象与事件，
是智能制造与机器感知领域的核心基础能力。随着深度学习在 3D 点云、RGB-D 多模态数据上的推进，异常检测任务
从传统单一传感器的阈值判别，走向融合几何结构、纹理与深度的高维空间建模，相关方法论（如自编码重建、
记忆库匹配、合成异常扩展、预训练特征微调等）层出不穷。

\begin{figure}[htbp]
\centering
\includegraphics[width=0.9\textwidth]{architecture.png}
\caption{SmartSurvey 系统分层架构示意图}
\label{fig:architecture}
\end{figure}

\subsection{文献量爆发带来的梳理痛点}
空间异常检测的子方向高度碎片化，文献以 PEG、阴极板检测、抓取质量评估、表面缺陷识别等多种形态散落在不同出版社中，
且在方法、指标与实验设置上缺乏统一口径。人工横向对比``哪篇方法在哪种传感器/时延约束下更优''需要反复翻阅原文，
效率极低且易遗漏关键证据。由此形成本系统的核心动机：以自动化手段把一组异构 PDF 转化为可复核、可对比、可直接成稿的结构化综述。

\subsection{系统目标与开源交付}
SmartSurvey 的目标是构建一个\textbf{自动化的学术文献综述系统}：输入任意数量、任意排版的 3D/RGB-D 空间异常检测论文 PDF，
输出为\textbf{空间视觉对比矩阵}与\textbf{中文 XeLaTeX 学术综述全文手稿}，支持 Markdown 实时预览、知识图谱可视化与
三件套导出（\texttt{survey\_draft.tex} / \texttt{matrix\_table.tex} / \texttt{references.bib}）。
项目已开源至 \url{https://github.com/C6-H14/SmartSurvey}，配备 \texttt{run\_windows.bat} 一键运行脚本与
Docker 容器镜像，可直接用于课程项目、科研选题与实验室调研。

% ============================================================
% 2. 系统架构与技术路线
% ============================================================
\section{系统架构与技术路线}

\subsection{分层架构与核心模块划分}
系统以 Python 3.10+ 与 Streamlit 为骨架，将职责拆分为可独立单测的 \texttt{core/} 分层模块，
形成``解析 → 提取 → 审查 → 生成''的流水线：

\begin{center}
\begin{tabularx}{\textwidth}{>{\bfseries}ll X}
\toprule
模块 & 职责 \\
\midrule
\texttt{core/pdf\_parser.py} & 核心章节识别 + 兜底页码切片（PyMuPDF） \\
\texttt{core/extractor.py} & 空间视觉矩阵提取 + 自愈重试（LLM） \\
\texttt{core/evidence.py} & 证据归一化与 containment 事实校验 \\
\texttt{core/schema.py} & 通用/领域双层 Schema 生成 \\
\texttt{core/credentials.py} & OS Keyring 凭据存储（无头环境降级） \\
\texttt{core/synthesis.py} & XeLaTeX 全文合成 + 语法栈扫描 + 物理编译自愈 \\
\texttt{core/templates.py} & LaTeX/BibTeX 模板与统一导言区（SSOT） \\
\texttt{core/graph.py} & PyVis 2D/3D 交互图形拓扑网络 \\
\texttt{core/pipeline.py} & 批处理流水编排（解析→提取→审查→生成） \\
\texttt{core/agent.py} & LLM 适配器（OpenAI 兼容，三级凭据回退） \\
\bottomrule
\end{tabularx}
\end{center}

在图形学技术选型上：PDF 解析使用 PyMuPDF（fitz），契合``章节识别 + 页码切片''混合策略；
知识图谱使用 \texttt{pyvis + networkx}——前者可将拓扑导出为交互 HTML，配合
\texttt{st.components} 嵌入网页渲染 2D/3D 图形网络。

\subsection{PyVis 交互图形可视化机制}
文献库（\texttt{data/vault\_100\_lab\_anomaly.json}）经 networkx 建模后，由 \texttt{core/graph.py}
布局为文献拓扑网络：节点表示单篇论文，边表示主题/方法关联，节点大小与颜色映射领域字段。
PyVis 输出独立 HTML，通过 \texttt{st.components.html} 内嵌于 Streamlit 面板，支持拖拽、缩放与悬浮高亮，
帮助研究者快速洞察文献聚类与主题演化。为适配无头容器，图谱采用懒加载与内存渲染策略。

\begin{figure}[htbp]
\centering
\includegraphics[width=0.9\textwidth]{p2.png}
\caption{Streamlit 主界面（侧边栏凭据表单 + 主面板矩阵与图谱预览）}
\label{fig:ui_main}
\end{figure}

\subsection{三级凭据降级与网络抗抖动设计}
面向公网部署与无头容器，系统设计了三级凭据降级链：
\[
\text{OS Keyring} \;\longrightarrow\; \text{.env 环境变量} \;\longrightarrow\; \text{st.session\_state 会话内存}
\]
本地桌面优先使用操作系统钥匙串，零明文落盘；无头 Linux 容器中捕获 \texttt{NoKeyringError}
平滑降级到会话内存，由外部注入环境变量/Secrets，保证公网零红屏崩溃。
.env 仅存放非机密的 \texttt{OPENAI\_API\_BASE} 与 \texttt{LLM\_MODEL\_NAME}，严禁写入真实密钥。

针对网络抖动，网关层实现指数退避重试：显式捕获 \texttt{InternalServerError} 与 \texttt{APITimeoutError}，
按 2s/4s/8s 三次退避并透传 \texttt{on\_retry} 钩子驱动前端提示，避免长文合成瞬间失败红屏。

% ============================================================
% 3. 核心算法与自动化合成
% ============================================================
\section{核心算法与自动化合成}

\subsection{空间视觉矩阵提取（双层 Schema）}
\begin{figure}[htbp]
\centering
\includegraphics[width=0.9\textwidth]{p0.png}
\caption{PyVis 2D/3D 交互图形拓扑网络渲染效果}
\label{fig:graph_render}
\end{figure}

面向空间异常检测这类异构领域，系统采用\textbf{通用字段 + 领域字段}的双层矩阵 Schema。
通用字段固定为 \texttt{title / authors / year / venue / research\_problem / method / innovation / limitation /
evidence\_page / evidence\_quote / confidence}；领域字段由用户输入的综述主题动态生成，例如工业空间异常检测方向
自动产出 \texttt{sensor}、\texttt{accuracy}、\texttt{latency}、\texttt{deployment\_scene}、
\texttt{decision\_mechanism} 等空间视觉指标，保证跨论文纵向对比的一致口径与横向对齐。字段缺失时标记为 \texttt{missing}，不以猜测填充。

\subsection{Inline Scope Binding 零丢纸事实校验}
为杜绝``证据套话''与``论文静默丢失''两类致命缺陷，系统施加两层硬约束。

\paragraph{证据包含校验（事实门控）} 每条局限与研究缺口必须绑定 \texttt{evidence\_page} 与
\texttt{evidence\_quote}（1--3 句英文原文摘录），后端执行严格包含校验：
$\text{evidence\_quote} \in \text{page\_text}[\text{evidence\_page}]$。
校验失败即判定为幻觉并触发自愈重试。

\paragraph{Inline Scope Binding（进程级物理绑定）} 在 \texttt{core/pipeline.py} 逐论文主循环内直接完成``提取—自愈—校验—就地降级''闭环：
每篇论文提取的行\textbf{立即}在当前循环做包含校验；
自愈 3 次失败的行就地降级，将 \texttt{limitation} 改写为 \texttt{"missing (unverified)"} 并强制保留。
由此保证最终产物行数等于输入 PDF 数，\textbf{100\% 论文不被静默丢弃}。

\subsection{双层模板 + XeLaTeX 全文动态排版与语法栈扫描器}
综述生成的机械核心是``单次 / 分章节多阶段''双层合成引擎与统一导言区（SSOT）架构，构成可动态排版的 XeLaTeX 物理排版引擎。

\begin{itemize}
  \item \textbf{统一导言区（SSOT）}：\texttt{\_build\_preamble()} 代码侧硬编码生成全部 LaTeX 导言区，
        剥离大模型生成导言区的权限；锁死 \texttt{margin=1.8cm}、\texttt{ctexart} 文档类与 \texttt{amsmath} 宏包，并置顶魔术注释。
  \item \textbf{单次合成（$\le 8000$ 字）}：一次 LLM 调用生成六章全文，之后执行\textbf{轻量语法栈扫描器}
        \texttt{validate\_latex\_syntax()}：逐字符校验 \texttt{\$} 奇偶、\texttt{\textbackslash begin/\textbackslash end}
        环境配对、花括号平衡，并检测中文右括号误写。
  \item \textbf{分章节多阶段合成（$>8000$ 字）}：将六章拆分为六次独立 LLM 调用，每次传入全量 rows 矩阵加上
        此前已生成章节的真实文本作为链式上下文，保证全文风格一致；结束符 \texttt{\textbackslash end\{document\}} 由代码自动追加。
  \item \textbf{物理 XeLaTeX 自愈}：\texttt{compile\_with\_xelatex()} 在临时目录中做非阻塞编译，
        从 .log 提取真实报错反馈给 LLM 触发自愈重试；编译器缺失或超时时静默降级为静态扫描。
  \item \textbf{跨论文辩证矛盾点分析}：在链式合成 Prompt 中显式要求大模型挖掘不同论文间的观点矛盾与技术权衡
        （如``论文 A 主张多模态融合能提升空间感知精度，而论文 B 证明多模态数据同步会导致边缘端延迟暴增 40\%''），
        将碎片化局限性升华为结构性矛盾。
\end{itemize}

在最终导出端，系统通过正则将 Markdown 粗体清洗为 \texttt{\textbackslash textbf\{\}}，
并实施以下多层物理清洗护栏：

\paragraph{RAG 标记清洗。}
\texttt{\_strip\_evidence\_page\_leaks()} 使用正则覆盖半角/全角圆括号、方括号与
CJK 角括号四种括号变体（含空白容错），\texttt{grep} 扫描零匹配，RAG 标记彻底物理清除。

\paragraph{数学符号规范化。}
\texttt{\_clean\_math\_operators()} 扫描数学模式将斜体 \texttt{or} 替换为 $\lor$ 符号。

\paragraph{参考文献自动护栏。}
\texttt{\_inject\_bibliography()} 在 \texttt{\textbackslash end\{document\}} 前自动注入含 [1]--[4] 四篇文献完整信息的
\texttt{thebibliography} 环境，幂等检查防重复注入。

\paragraph{矩阵表格化。}
第四节从 description 列表重构为 \texttt{tabularx} 五列三线表（文献\&年份/异常范式/模态输入/关键指标/局限性）。

% ============================================================
% 4. 实验测试与运行效果分析
% ============================================================
\section{实验测试与运行效果分析}

\subsection{测试环境}
系统在异构目标形态下完成部署与验证，覆盖本地桌面与容器两类运行环境：
\begin{itemize}
  \item \textbf{Windows 本地}：\texttt{run\_windows.bat} 一键启动，自动校验 Python 3.10+、创建
        \texttt{.venv} 虚拟环境并安装依赖；显式锁定 \texttt{.venv\textbackslash Scripts\textbackslash python.exe} 路径消除环境污染。
  \item \textbf{Docker 容器}：基于 \texttt{python:3.11-slim}，暴露 8501 端口；
        无头容器中图谱按内存渲染，密钥经环境变量注入。
  \item \textbf{Streamlit Cloud}：针对 0.5 vCPU 与 60s 超时，引入 \texttt{@st.cache\_data} 缓存、
        图谱懒加载与 120s 客户端 Timeout。
\end{itemize}

\begin{figure}[htbp]
\centering
\includegraphics[width=0.9\textwidth]{p5.png}
\caption{终端运行 \texttt{python -m pytest tests -v} 全绿（126 passed）}
\label{fig:terminal_pytest}
\end{figure}

\subsection{126 项自动化测试与覆盖率}
系统以``测试即规约''贯穿开发，共 \textbf{126 项自动化单元测试全部通过（100\% PASS）}，覆盖：
\textbf{解析层}（章节识别失败返回 \texttt{missing}；非标准 PDF 保留 page slices）；
\textbf{证据校验层}（\texttt{evidence\_quote in page\_text} 门控，零丢纸保证 N 篇 PDF 恰好 N 行）；
\textbf{凭据安全层}（API Key 不出现于日志/导出文件；无头降级路径为会话内存）；
\textbf{语法扫描层}（\texttt{validate\_latex\_syntax()} 覆盖未闭合 \texttt{\$}、环境不匹配、花括号不平衡、CJK 误写）；
\textbf{网络抗抖动层}（端到端冒烟测试 \texttt{test\_e2e\_smoke.py} 拦截参数签名脱节等集成性缺陷）。

\subsection{运行效果与产出生手稿深度分析}

在流式提取 4 篇空间异常检测文献并生成综述的基准流程中，系统依次执行解析 → 提取 → 证据审查 → 动态排版，
产出含六章节的中文 XeLaTeX 手稿与空间视觉对比矩阵。以下从真实文献提炼、形式化数学合成、辩证批判三个维度，展示系统实际产出的核心学术价值。

\subsubsection{真实文献与范式提炼}

系统自动从输入 PDF 中提取并结构化整理了 4 篇代表性工作，覆盖空间异常检测的四个核心技术范式：

\noindent\textbf{Soudani 等：单模态 YOLO 实时监控。}
以 YOLO 单阶段检测器在 RGB 帧上直接输出空间异常边界框与置信度，30+ FPS 满足工业实时需求。
系统提取其核心创新为``将空间异常检测转化为单帧目标检测''，证据摘录精确绑定原文第 4 页；
同时识别其局限：仅依赖 RGB 模态，低光照/遮挡下漏检率显著上升，且缺乏时序连续性建模。

\noindent\textbf{Costanzino 等：跨模态特征映射。}
提出将 RGB 与深度/热红外映射至共享特征空间，经跨模态注意力融合提升感知精度。
系统提取传感器类型为 RGB-D-T，准确率指标 mAP@0.5 = 89.3\%。
该方法的形式化贡献——跨模态特征差异度量——被系统自动识别为领域关键公式（见下文）。
其局限在于多模态数据同步要求传感器阵列严格标定，工业现场振动与温漂下误差累积显著。

\noindent\textbf{Bergmann 等：双分支逻辑异常检测。}
以 MVTec 3D-AD 为基础，几何分支处理点云/深度图表面缺陷，纹理分支处理 RGB 外观异常，
两分支经逻辑与门融合。系统识别其部署场景为``工业质检流水线''，决策机制为``双分支逻辑与 + 阈值判决''。

\noindent\textbf{Iodice 等：行为树自适应决策。}
引入行为树替代固定阈值级联，根据异常类型（位置偏差/姿态异常/表面缺陷）与置信度动态选择响应路径
（忽略/记录/报警/停机），将平均决策延迟从 420 ms 降至 180 ms，映射至领域字段 \texttt{decision\_mechanism}。

\subsubsection{形式化数学合成}

系统双层模板合成引擎成功识别并渲染了跨论文的形式化数学关系，经语法栈扫描器校验后嵌入综述手稿：

\paragraph{跨模态特征差异度量（Costanzino 等）。}
RGB 与深度/热红外模态的特征对齐误差形式化为：
\begin{equation}
s(x) = \|f_B(x) - M(f_A(x))\|_2
\label{eq:crossmodal}
\end{equation}
其中 $f_A$ 为 RGB 支路特征提取器，$f_B$ 为深度/热红外支路特征提取器，$M$ 为跨模态映射网络。
$s(x)$ 越大表明跨模态融合越不可靠。系统自动将式 (\ref{eq:crossmodal}) 与 Costanzino 的 \texttt{limitation}
关联，指出 $s(x) > \tau$ 时多模态融合退化为单模态 YOLO 基线，精度回退至 Soudani 水平。

\paragraph{可微软约束损失函数（Bergmann 等）。}
从双分支架构中抽象出带逻辑约束的联合优化目标：
\begin{equation}
\mathcal{L}_{\text{total}} = \mathcal{L}_{\text{detection}} + \lambda \sum_{j=1}^{M} (1 - g_j(z))^2
\label{eq:softconstraint}
\end{equation}
其中 $\mathcal{L}_{\text{detection}}$ 为几何与纹理分支联合检测损失，
$g_j(z) \in [0,1]$ 为第 $j$ 条逻辑约束的满足度（如``几何异常 + 纹理正常 → 非缺陷''），
$\lambda$ 为约束强度。该式将不可微的逻辑与门松弛为可微软约束，使双分支网络可端到端训练。
系统注明此公式为从 Bergmann 方法论中归纳的形式化抽象，体现了综述系统的方法论提炼能力。

\subsubsection{深度辩证批判}

系统在链式合成的``跨论文辩证矛盾点分析''阶段，成功梳理出两大技术张力：

\paragraph{张力一：多模态精度 vs.\ 系统脆弱性。}
Costanzino 等主张多模态融合将 mAP 从单模态 78.1\% 提升至 89.3\%，但强依赖传感器完备性与标定精度；
Soudani 等的单模态方案精度虽低，在传感器退化时仍可维持基线运行。系统综述明确指出：
``多模态方法的精度增益建立在对传感器阵列完备性的强假设之上；一旦模态缺失，性能退化至远低于专用单模态方案。''
这一分析将碎片化 \texttt{limitation} 升华为结构性方法论张力。

\paragraph{张力二：实时决策 vs.\ 高精度推理延迟。}
Iodice 等行为树方案以 180 ms 决策延迟实现工业实时响应，但基于启发式规则，缺乏深度推理；
Bergmann 等双分支架构虽可精确判别几何 + 纹理双重异常（准确率 94.2\%），推理延迟约 1.2 s/帧，
难以满足高速产线（$>$1 件/秒）。系统将此归纳为``精度--延迟帕累托前沿''，
并自动建议知识蒸馏与模型量化作为弥合路径——该建议为系统基于跨论文矛盾综合生成的原创判断，
展现了综述引擎从``对比''到``洞察''的认知闭环。

综上，系统产出综述不仅实现结构完整性（3000--5000 字、6 章标准结构）
与事实校验（100\% 证据绑定），更在文献提炼、数学合成与辩证批判三层面展现了接近人工深度的综合分析能力。

\subsubsection{排版质量护栏升级验证}

审阅专家反馈驱动的三项护栏升级已完成合成验证：

\paragraph{RAG 标记零泄漏。}
\texttt{\_strip\_evidence\_page\_leaks()} 升级为四种括号变体通用清洗后，
\texttt{grep evidence\_page survey\_draft.tex} 零匹配——RAG 标记彻底物理清除。

\paragraph{参考文献自动补全。}
\texttt{\_inject\_bibliography()} 在 \texttt{\textbackslash end\{document\}} 前自动注入
[1]--[4] 四篇完整出版信息的 \texttt{thebibliography} 环境，编译后 PDF 参考文献规范对齐。

\paragraph{公式与表格规范化。}
\texttt{\_clean\_math\_operators()} 将数学模式中 \texttt{or} 替换为 $\lor$；
学术对比矩阵重构为 \texttt{tabularx} 五列三线表，编译零 Overfull 警告。

% ============================================================
% 5. 总结与未来展望
% ============================================================
\section{总结与未来展望}

\subsection{系统贡献总结}
SmartSurvey 面向空间异常检测领域完成了从``批量 3D/RGB-D 论文 PDF''到``中文 XeLaTeX 综述手稿''的自动化闭环，
核心贡献归纳如下：
\begin{itemize}
  \item 以``核心章节识别 + 兜底页码切片''的混合解析与双层空间视觉矩阵 Schema，实现对异构空间异常检测文献的泛化梳理。
  \item 以``页码 + 原文''证据包含校验与 Inline Scope Binding，实现零丢纸、零套话的严谨事实门控。
  \item 以 PyVis 构建 2D/3D 交互图形拓扑网络，将文献主题关联直观可视化。
  \item 以双层模板 XeLaTeX 物理图形排版引擎（单次/多阶段 + 统一导言区 SSOT + 语法栈扫描 + 物理自愈 + 跨论文辩证矛盾点分析）实现``贴入 Overleaf 即可编译''的交付目标，并在综述中自动挖掘论文间的观点冲突与技术权衡。
  \item 以三级凭据降级与网关指数退避重试，保障本地桌面、Docker 容器与公网云三种形态下的安全与稳健运行。
\end{itemize}

\subsection{未来展望}
\begin{itemize}
  \item \textbf{边缘智联计算平台部署}：将综述系统迁移至边缘智联计算平台（如联宝 EA-B310），
        结合本地 LLM 推理与轻量化图谱渲染，实现离线化文献智能服务。
  \item \textbf{机械臂空间视觉部署}：将综述中提取的 \texttt{sensor / latency / deployment\_scene}
        指标对齐到实际机械臂视觉标定与在线推理链路，探索``综述洞见 → 工程落地''的正向闭环。
  \item \textbf{自动文献获取}：接入学术搜索引擎自动抓取全文并落地 Zotero 归档，减少人工整理成本。
  \item \textbf{更深度的批判性评述}：引入置信度加权与自动反事实推理，提升综述的批判性学术深度。
\end{itemize}

% ============================================================
% 致谢
% ============================================================
\section*{致谢}
感谢中国科学技术大学《计算机图形学前沿》课程主讲老师刘利刚教授及各位助教的精彩课程。
感谢翟晓雅老师的细心指导及指出的``可验证最小闭环''等工程要点。
感谢陈现敏老师提供大学生研究计划（大研）机会并协助选定空间异常检测研究方向。
感谢软件工程学院《AI4SE》课程传授的 Superpowers 技能框架与 Spec-Driven 开发范式。
感谢测试阶段提供真实论文样例并反馈意见的各位同学。

% ============================================================
% 参考文献
% ============================================================
\begin{thebibliography}{9}

\bibitem{pymupdf}
Artifex. \textit{PyMuPDF Documentation}[EB/OL]. \\ \url{https://pymupdf.readthedocs.io/}, 2026.

\bibitem{pyvis}
West Health Institute. \textit{PyVis: Interactive Network Visualization}[EB/OL]. \\ \url{https://pyvis.readthedocs.io/}, 2026.

\bibitem{streamlit}
Streamlit Inc. \textit{Streamlit Data App Framework}[EB/OL]. \\ \url{https://docs.streamlit.io/}, 2026.

\bibitem{mvtec3d}
Bergmann P, Batzner K, Fauser M, et al. \textit{The MVTec 3D-AD Dataset for Unsupervised 3D Anomaly Detection and Localization}[C]. VISAPP, 2022.

\bibitem{langchain}
LangChain Project. \textit{LangChain \& OpenAI Integration Docs}[EB/OL]. \\ \url{https://python.langchain.com/}, 2026.

\end{thebibliography}

\end{document}